\documentclass[livre]{udes-genie-these}

\usepackage[T1]{fontenc}
\usepackage{booktabs}

\ConfigurationDocument
{
  langue = francais,
  type = memoire,
  programme = electrique,
  lignes,
  liste-figures,
  liste-tableaux,
  fichier-resume-francais = {resume-francais},
  %fichier-resume-anglais = {resume-anglais},
  fichier-remerciements = {merci},
  fichier-lexique = {lexique},
  fichier-symboles = {symboles},
  fichier-acronymes = {acronymes},
  fichiers-references = {references},
  %auto-bib,
  %style-references = plain,
  %legendes-min,
}

\TitreFrancais{CECI EST LE TITRE FRANÇAIS\\DU DOCUMENT}
%\TitreAnglais{DOCUMENT ENGLISH TITLE}
\Auteur{Prénom}{NOM}
\Date{Septembre}{2022}
\Dedicace{À mon chien préféré, Pitou, et à mon chat préféré, Minou.}
\MotsClesFrancais{3 à 8 mots-clés français}
%\MotsClesAnglais{3 to 8 english keywords}

\Directeur{Prénom1}{NOM1}
\Codirectrice{Prénom2}{NOM1}
\Codirecteur{Prénom3}{NOM3}
\Evaluatrice{Prénom4}{NOM4}
\Evaluateur{Prénom5}{NOM5}

% Le contenu du document doit respecter le protocole de rédaction.
% Le contenu fourni ici n'est qu'un exemple afin d'illustrer l'utilisation de la classe.

\begin{document}

\chapter{INTRODUCTION}

Ce document est uniquement un exemple d'utilisation de la classe \LaTeX\ pour la rédaction de documents conforme au \textit{Protocole de rédaction aux études supérieures} de la Faculté de génie de l'Université de Sherbrooke.  Le contenu de cet exemple n'est pas représentatif de ce que devrait contenir un vrai document. Pour connaître le contenu éventuel d'un vrai document, il faut se référer au protocole.

Les références bibliographiques qui suivent ne sont que des exemples. Il faut se référer au protocole pour plus de détails à ce sujet afin de s'assurer que les références sont complètes et qu'elles sont conformes aux exigences.

\begin{description}
  \item[Article de périodique:]
  il y a quatre exemples \cite{burke1995,pattamatta2009,roy2009,vasu2009}.

  \item[Articles d'actes de conférence:]
  il y a un exemple court~\cite{prevenslik2009}
  et deux exemples longs~\cite{coronado2005,yeung2006}.

  \item[Livre:]
  il y a un exemple court~\cite{russell2003}
  et trois exemples longs~\cite{kececioglu1991,knuth1998-num,meriam2010}.

  \item[Livre, une partie:]
  il y a deux exemples courts~\cite{cormen2001,gaddis2009}
  et trois exemples longs~\cite{kececioglu1991-partie,knuth1998-sort,meriam2010-dist}.

  \item[Essai, mémoire et thèse:]
  il y a un exemple de chacun \cite{belanger2009,chafei2009,francoeur2009}.

  \item[Rapport technique:]
  il y a un exemple court~\cite{belley2008}
  et un exemple long~\cite{kernighan1991}.

\end{description}

\chapter{CONCEPTION}

\chapter{RÉALISATION}

\chapter{TESTS}

\chapter{ANALYSE}

\chapter{CONCLUSION}

\appendix

\chapter{DONNÉES}

\end{document}